\documentclass[11pt]{article}
\usepackage[margin=1in]{geometry}
\usepackage{amsmath,amssymb}
\usepackage{booktabs}
\usepackage{hyperref}
\usepackage{listings}
\usepackage{xcolor}
\lstset{basicstyle=\ttfamily\small,breaklines=true,frame=single,columns=fullflexible}
\title{Independent Replication of \\ \emph{Quantum Algorithms} (Mosca, arXiv:0808.0369) \\ QC-200 Wave, Survey-Mode Spot-Check}
\author{Ollie (subagent replication) --- 2026-07-05}
\date{}
\begin{document}
\maketitle

\begin{abstract}
Michele Mosca's ``Quantum Algorithms'' (\href{https://arxiv.org/abs/0808.0369}{arXiv:0808.0369v1}, 2008) is a broad expository \emph{survey} of the quantum-algorithms landscape as of mid-2008 (Shor factoring / order-finding / discrete log / Abelian HSP, Grover / amplitude amplification, quantum simulation, quantum walks, adiabatic and topological algorithms, quantum-tasks algorithms). Because a survey has no single novel empirical claim to falsify, we perform a \textbf{spot-check reproduction}: we implement, in Qiskit statevector simulation, two representative algorithms whose quantitative behaviour Mosca describes, and we check whether the numbers he cites (or that follow immediately from the formulas he quotes) come out. Both spot-checks match theory exactly to machine precision:
(i) 3-qubit Grover search on $N{=}8$ with one marked item hits $P(|w\rangle){=}0.94531$ after $k{=}\lfloor(\pi/4)\sqrt{N}\rfloor{=}2$ iterations, matching $\sin^2((2k{+}1)\theta)$ with $\theta{=}\arcsin(1/\sqrt{8})$;
(ii) Shor order-finding for $a{=}7$, $N{=}15$ with an 8-qubit counting register produces four peaks of probability exactly $1/r{=}0.25000$ at $k\in\{0,64,128,192\}{=}\{k\cdot 2^m/r\}$, and continued-fraction recovery on the top outcomes yields $r{=}4$ as expected. \textbf{Verdict: SPOT-CHECK} --- both sub-algorithms whose formulas we exercised match the paper-quoted numbers exactly.
\end{abstract}

\section{Paper Summary}
Mosca 2008 is a $\sim$61-page survey (bibliography [1]--[188]) covering the state of quantum algorithms circa 2008 and reprinted in the \emph{Encyclopedia of Complexity and System Science}. It is structured as:
\begin{enumerate}\itemsep=0pt
  \item Definition of the Subject; Introduction and Overview
  \item Early quantum algorithms (Deutsch, Deutsch--Jozsa, Bernstein--Vazirani, Simon)
  \item Factoring, discrete logarithms, Abelian Hidden Subgroup Problem (Shor's algorithms; Kitaev's phase-estimation formulation)
  \item Algorithms based on amplitude amplification (Grover, Brassard--H\o yer--Mosca--Tapp)
  \item Simulation of quantum-mechanical systems
  \item Generalisations of the Abelian HSP (non-Abelian HSP, Pell's equation, hidden shift, etc.)
  \item Quantum-walk algorithms (element distinctness, triangle finding, formula evaluation)
  \item Adiabatic algorithms
  \item Topological algorithms (Jones polynomial approximation)
  \item Quantum algorithms for quantum tasks (state tomography, spectrum estimation)
\end{enumerate}
As a survey the paper contains \emph{no new empirical result of the author's own}. Its testable content is the collection of quantitative bounds and success probabilities it \emph{quotes} from the primary literature (e.g. $O(\sqrt{N})$ Grover queries; $\Omega(\sqrt{N})$ lower bound; $O((\log N)^3)$ Shor query complexity; the specific formula for the number of Grover iterations that maximises the success probability). We therefore treat this as a spot-check: pick two of the algorithms Mosca describes in detail, implement them at a small faithful instance, and verify the reproduced numbers agree with the closed-form theory Mosca (implicitly or explicitly) cites.

\section{Claims Table}
\begin{center}
\begin{tabular}{@{}p{0.7cm}p{6.8cm}p{2.2cm}p{2.2cm}p{2.4cm}@{}}
\toprule
\# & Claim (paraphrased from Mosca 2008) & Type & Testable in Qiskit @ small $N$? & Tested here? \\
\midrule
C1 & Grover's algorithm finds a marked item in an unstructured search space of size $N$ in $O(\sqrt{N})$ oracle calls; the optimal number of iterations for $M{=}1$ marked item is $k^*=\lfloor(\pi/4)\sqrt{N/M}\rfloor$, achieving success probability $\sin^2((2k^*{+}1)\theta)$ with $\theta=\arcsin(\sqrt{M/N})$. & Quantitative & Yes ($n{=}3$) & \textbf{YES} \\
C2 & Shor's factoring algorithm reduces to order-finding: for $a$ coprime to $N$, the QFT on a counting register after applying controlled-$U^{2^j}$ (where $U|y\rangle{=}|ay\bmod N\rangle$) yields, upon inverse QFT and measurement, values $k$ concentrated at multiples of $2^m/r$, where $r$ is the order of $a\bmod N$. Continued-fraction expansion of $k/2^m$ recovers $r$. & Quantitative & Yes ($N{=}15$, $a{=}7$, $m{=}8$) & \textbf{YES} \\
C3 & Amplitude amplification generalises Grover: given a heuristic algorithm $\mathcal A$ succeeding with probability $p$, one recovers a solution in $O(1/\sqrt p)$ calls to $\mathcal A,\mathcal A^{-1},U_f$. & Quantitative & Yes (algorithmic corollary of C1) & Implicit via C1 \\
C4 & Simulation of quantum systems on quantum hardware requires only $\mathrm{poly}(n)$ gates for local Hamiltonians. & Structural/complexity & Not empirical & No (out of scope of a survey spot-check) \\
C5 & Quantum walks give speed-up over classical for element distinctness ($O(N^{2/3})$) and triangle finding. & Complexity & Small-$N$ demos possible but $\gg$ scope & No \\
C6--C$k$ & Non-Abelian HSP / adiabatic / topological / tomography sections. & Complexity/structural & Mostly out of scope for a survey spot-check & No \\
\bottomrule
\end{tabular}
\end{center}

We exercise \textbf{C1} (Grover, quantitatively) and \textbf{C2} (order-finding, quantitatively). C3 is implied by C1's exact match to the amplitude-amplification formula. C4--C6 are not spot-checked here; a full replication of a survey would need $\sim$10 sub-projects (roughly the QC-100 wave itself).

\section{Method}
\subsection{Environment and tool versions}
\begin{itemize}\itemsep=0pt
  \item Host: CherryRd (macOS Darwin 25.3.0, Python 3.13 in a fresh venv \verb|.venv/|).
  \item \textbf{Qiskit 2.5.0} (via \verb|pip install qiskit qiskit-aer numpy|). Statevector-only, no shots, no noise.
  \item Poppler \verb|pdftotext| for PDF text extraction fallbacks (Marker / Nougat were not installed on this host; see \verb|extraction/marker.md| and \verb|extraction/nougat.mmd| for the fallback stanza).
  \item No LLM calls needed for the reproduction itself; had they been needed, Argo (localhost:44497 key=stevens) would have been used per the standing free-endpoints rule.
\end{itemize}

\subsection{Spot-check 1 --- Grover on $N{=}8$, $M{=}1$}
Script: \verb|report/evidence/grover_N8.py|.
Circuit: 3 qubits. Initialise to $|s\rangle{=}H^{\otimes 3}|000\rangle$. Apply the Grover iterate $G{=}D\cdot O_w$ twice ($k^*{=}\lfloor(\pi/4)\sqrt 8\rfloor{=}2$), where $O_w$ is the phase-flip oracle for the marked state $|w\rangle{=}|101\rangle{=}|5\rangle$, and $D{=}H^{\otimes 3}S_0 H^{\otimes 3}$ with $S_0{=}2|0\rangle\langle 0|{-}I$ (implemented as X-conjugated multi-controlled-Z per the standard construction). Extract the exact statevector $|\psi\rangle$ via \verb|Statevector.from_instruction|, read off $|\langle w|\psi\rangle|^2$.

\subsection{Spot-check 2 --- Order-finding for $a{=}7$, $N{=}15$}
Script: \verb|report/evidence/orderfinding_a7_N15.py|.
Circuit: $m{=}8$ counting qubits $+$ $n_\text{work}{=}\lceil\log_2 N\rceil{=}4$ work qubits. Superpose the counting register with $H^{\otimes m}$; initialise the work register to $|1\rangle$. For each counting-bit $j{=}0,\dots,m{-}1$, append a controlled-$U^{2^j}$, where $U|y\rangle=|ay\bmod N\rangle$ for $y<N$ and $U|y\rangle=|y\rangle$ for $y\in\{N,\dots,2^{n_\text{work}}{-}1\}$; we build $U^{2^j}$ directly as the permutation matrix induced by multiplication by $a^{2^j}\bmod N$ and wrap it with \verb|UnitaryGate(...).control(1)|. Apply the inverse QFT on the counting register. Extract the full statevector, marginalise over the work register, list the highest-probability counting outcomes, and run continued-fraction expansion of $k/2^m$ to recover candidate orders.

\section{Results vs.\ Paper}
\begin{center}
\begin{tabular}{@{}p{0.5cm}p{6.5cm}p{3.5cm}p{3.5cm}p{1.2cm}@{}}
\toprule
\# & Quantity & Paper / theory prediction & Our simulation & Match \\
\midrule
C1 & Grover success probability for $N{=}8$, $M{=}1$, $k{=}2$ & $\sin^2(5\theta)$ with $\theta{=}\arcsin(1/\sqrt 8)$ = \textbf{0.94531} & \textbf{0.94531} (agreement to $10^{-6}$) & \checkmark \\
C1 & Amplitude on $|w\rangle$ & $\sin(5\theta)=0.97227$ & $0.97227{+}0j$ & \checkmark \\
C1 & Gain factor over uniform $1/N{=}0.125$ & $\times$ 7.5625 & $\times$ 7.5625 & \checkmark \\
C2 & Number of QFT peaks & $r{=}4$ (classical: $\mathrm{ord}_{15}(7){=}4$) & 4 peaks (all others $<10^{-30}$) & \checkmark \\
C2 & Peak positions $k=\lfloor s\cdot 2^m/r\rfloor$ for $s{=}0,1,2,3$ & $\{0,64,128,192\}$ & $\{0,64,128,192\}$ & \checkmark \\
C2 & Probability at each peak & $1/r{=}0.25000$ & $0.25000\pm 10^{-14}$ & \checkmark \\
C2 & Continued-fraction recovery of $r$ from top peaks & $r=4$ & candidates $\{1,2,4\}$ (with $4$ present) & \checkmark \\
\bottomrule
\end{tabular}
\end{center}

\subsection{Grover raw output (Statevector, $N{=}8$)}
\begin{lstlisting}
n=3 N=8 marked=5 (bin 101) M=1 k_opt=2
P(|101>) after k=2 Grover iterations = 0.945312
amplitude on |101> = (0.9722718241315007+0j)
full probability distribution:
  |000> P=0.007812  |001> P=0.007812  |010> P=0.007812  |011> P=0.007812
  |100> P=0.007812  |101> P=0.945312  |110> P=0.007812  |111> P=0.007812
Theory: theta = arcsin(sqrt(1/8)) = 0.361367
Theory: amplitude on marked = sin(5*theta) = 0.972272
Theory: P(marked)          = 0.945312
MATCH (sim vs theory, rel_tol=1e-6): True
\end{lstlisting}
Every unmarked basis state carries residual probability $(1{-}0.9453)/7\approx 0.00781$, exactly the $\cos^2((2k{+}1)\theta)/(N{-}1)$ prediction.

\subsection{Order-finding raw output ($a{=}7$, $N{=}15$, $m{=}8$)}
\begin{lstlisting}
Classical order r of a=7 mod N=15 is r=4
a^(2^j) mod N for j=0..7: [7, 4, 1, 1, 1, 1, 1, 1]

Top counting-register outcomes:
  k=64   P=0.250000   k/2^m=0.250000   ~ 1/4
  k=0    P=0.250000   k/2^m=0.000000   ~ 0
  k=192  P=0.250000   k/2^m=0.750000   ~ 3/4
  k=128  P=0.250000   k/2^m=0.500000   ~ 1/2

Expected peak positions (k*2^m/r) for r=4: [0, 64, 128, 192]
Sum of prob at expected peaks: 1.000000
Recovered candidate orders (denominators of continued-fraction rounds): [1, 2, 4]
Overall MATCH: True
\end{lstlisting}
Because $r{=}4$ divides $2^m{=}256$ exactly, the peaks are ``perfect'' (probability exactly $1/r$ each and summing to 1). For values of $r$ that do not divide $2^m$ evenly (e.g.\ $a{=}2,N{=}21,r{=}6$), the peaks broaden but continued fractions still recover $r$ with high probability; we did not perform that harder case since Mosca's exposition only asserts the ``multiples of $1/r$'' behaviour, which the $r{=}4$ case exercises cleanly.

\section{Verdict}
\textbf{SPOT-CHECK.} Both Qiskit reproductions match the closed-form predictions Mosca quotes to machine precision (Grover $0.94531$ vs $\sin^2(5\theta){=}0.94531$; order-finding four peaks at $\{0,64,128,192\}$ each of probability exactly $1/r{=}0.25$, continued fractions recovers $r{=}4$). Because Mosca 2008 is a survey with no unique novel empirical claim of the author's own, ``REPLICATED'' is not the appropriate verdict --- \textbf{SPOT-CHECK} is the correct verbal category per the wave brief. We did not exercise C4--C6 (simulation, quantum walks, adiabatic, topological, non-Abelian HSP); a full survey replication would require running the QC-100 or QC-200 waves in their entirety, which is out of scope for a single-paper subagent.

\section{Artifacts}
See \verb|report/artifacts_summary.md| for the full inventory. Code lives in \verb|report/evidence/*.py|; raw stdout logs in \verb|report/evidence/*.log|; JSON results in \verb|report/evidence/*_result.json|; PDF at \verb|paper.pdf|; Marker / Nougat fallback extractions in \verb|extraction/|.

\section*{Open Questions}
\begin{itemize}
  \item \textbf{Q1} --- For order-finding with $r\nmid 2^m$ (the ``non-trivial'' case Mosca alludes to but does not spell out numerically), what is the smallest $m$ such that one measurement suffices for continued-fraction recovery of $r$ with $\ge 90\%$ probability, as a function of $r$ and $N$? Our $r{=}4$ case is misleadingly clean because $4\mid 256$.
  \item \textbf{Q2} --- Mosca's ``$O(\sqrt N)$'' Grover claim collapses to $\lfloor(\pi/4)\sqrt N\rfloor$ for known $M$; for \emph{unknown} $M$ (which the survey mentions), what is the worst-case ratio of the exact optimal iteration count to the count chosen by the standard exponential-doubling BBHT strategy over the range $M{\in}[1,N/4]$, and does that ratio approach $\pi/8$ as $N\to\infty$?
  \item \textbf{Q3} --- Our controlled-$U^{2^j}$ implementation used a raw $16{\times}16$ permutation-matrix \verb|UnitaryGate|; how does the transpiled two-qubit-gate count scale if we instead synthesise controlled-multiplication-mod-$N$ from CCX + adders (Beauregard-style)? Mosca cites $O((\log N)^3)$; is there a small-$N$ regime where the raw permutation is actually cheaper on real hardware after Qiskit's transpiler pass?
  \item \textbf{Q4} --- The survey groups ``simulation of quantum-mechanical systems'' (Section 6) under quantum-algorithms speed-ups, but the 2008 exposition predates Low--Chuang qubitization and the linear-combination-of-unitaries family. Which of the specific complexity claims Mosca quotes (e.g.\ Trotter step count for a $k$-local Hamiltonian) have been tightened in the intervening years, and is there a modern survey we should read alongside 0808.0369 to close the gap?
  \item \textbf{Q5} --- Both spot-checks were exact-precision \emph{statevector} runs; the numbers reproduce to $\sim 10^{-14}$. Under a realistic depolarising-plus-readout noise model at fixed $p_1$ and $p_2$ (say $10^{-3}$), at what noise level does the Grover $k{=}2$ success probability drop below the classical $1{/}N{=}0.125$ baseline? Answering this would let us plot a ``useful-Grover boundary'' as a function of $N$ and noise, which would be a concrete pedagogical companion to Mosca's exposition.
\end{itemize}

\end{document}
